Suppose first that $Y$ is not a line. Choose affine coordinates with $P=(0,0)$ and write its equation as $f=ax+by+$ terms of degree at least $2$, where $(a,b)\ne(0,0)$. Restricting to a line through $P$ gives intersection multiplicity $1$ unless its direction annihilates $ax+by$. For the unique line $ax+by=0$, the restriction has order greater than $1$. A line not passing through $P$ has multiplicity $0$ there. This proves uniqueness.

For a homogeneous equation $F$, Euler's identity identifies this line with
\[F_{x_0}(P)x_0+F_{x_1}(P)x_1+F_{x_2}(P)x_2=0.\]
The partial derivatives are homogeneous of the same degree and do not vanish simultaneously on $\operatorname{Reg}Y$. Hence $P\mapsto(F_{x_0}(P):F_{x_1}(P):F_{x_2}(P))$ is a morphism to $(\mathbb{P}^2)^*$.

If $Y$ is a line, its tangent line is $Y$ itself and the morphism is constant. No distinct line has intersection multiplicity greater than $1$ with $Y$.
